\documentclass{article}
\PassOptionsToPackage{round,authoryear}{natbib}
\usepackage[dblblindworkshop]{neurips_2026}
\workshoptitle{Self-Evolving Diversity-Driven Search for Robust AI Systems (EvoRobust 2026)}

\usepackage[utf8]{inputenc}
\usepackage[T1]{fontenc}
\usepackage{hyperref}
\usepackage{url}
\usepackage{booktabs}
\usepackage{amsmath,amssymb}
\usepackage{amsfonts}
\usepackage{microtype}
\usepackage{graphicx}

\newtheorem{lemma}{Lemma}

\title{Diversity You Cannot See From the Mean: Calibrated Detection of \\
Latent Failure-Mode Structure in Open-Model Ecosystems}

\author{
  Anonymous Author(s) \\
  Affiliation \\
  Address \\
  \texttt{email}
}

\begin{document}

\maketitle

\begin{abstract}
Standard mean-level diversity statistics for model panels --- double-fault, pairwise disagreement,
observed agreement --- are computed the same way whether the panel is meant to be diverse or not:
count how often members fail, or answer, alike. We show this class of statistic is provably
degenerate for the failure case: mean pairwise co-failure is a deterministic function of item
difficulty alone and carries no information about shared behavior once that is accounted for. On
$1{,}228$--$1{,}362$ open language models across five benchmarks, the mean-level ``excess''
co-failure a naive computation reports is identically zero under the correct null, yet a genuine,
calibration-only-visible residual correlation of $2.8$--$11.8\times$ its null survives, concentrated
in $4$--$12$ latent behavioral dimensions depending on benchmark. An uncalibrated version of the
same summary statistic reads $1.5$--$3.6$ on every benchmark before conditioning --- easily
misread as near-total diversity --- and only calibration against an exact, fit-free reference
distribution reveals the real structure underneath. We argue this measurement-validity problem
transfers directly to diversity-aware and continuously evolving safety benchmarks wherever coverage
is scored from pairwise (dis)agreement rates, and we describe the calibration as a drop-in
diagnostic rather than a replacement for search-based diversity methods.
\end{abstract}

\section{The measurement problem}

A recurring move in both classifier-ensemble diversity and, increasingly, in diversity- and
coverage-aware safety evaluation is to summarize a panel's behavioral spread from pairwise
(dis)agreement: how often do two members fail, or answer, the same way. This workshop's own scope
lists ``diversity and coverage metrics for safety evaluation'' and ``behavioral repertoires of
failures'' as central objects of study. We show that the most common member of this family --- mean
pairwise co-failure, one of the ten classifier-diversity measures catalogued by
\citet{kuncheva2003} under the name double-fault --- has a property that makes it a poor diversity
signal on its own: it is a deterministic function of item-level difficulty alone.

\begin{lemma}[Mean co-failure depends only on item margins]
Let $F_{im}=1$ iff panel member $i$ fails item $m$, with $c_m$ the number of members failing item
$m$. Then $O = \frac{1}{MN(N-1)}\sum_m(c_m^2-c_m)$, a function of the column counts $\{c_m\}$
alone.
\end{lemma}
\noindent Consequently, any randomization of which member fails which item that preserves each
item's failure count leaves $O$ exactly invariant, with zero variance. A benchmark with a few very
hard items will show high mean co-failure even if every member fails independently given the
item; a benchmark with uniformly moderate items will show low co-failure even if members are
identical clones. The statistic conflates \emph{item difficulty} with \emph{shared behavior}, and
no amount of more data resolves this — it is an identity, not a sampling artifact.

\section{A calibrated diversity signal, at scale}

The fix is to condition on the item margins exactly rather than compare against an uninformative
uniform baseline. The row and column margins of a binary outcome matrix are the jointly sufficient
statistics of the Rasch model, so the uniform distribution over margin-matched matrices --- sampled
here with curveball \citep{strona2014}, verified against direct enumeration on small fibres --- is
the unique exact, fit-free null that removes item difficulty without fitting anything. We apply
this to the archived Open LLM Leaderboard: $1{,}228$--$1{,}362$ open models per benchmark, five
benchmarks, harvested from result files at zero inference cost.

\begin{table}[t]
\centering
\small
\begin{tabular}{lrrrrr}
\toprule
Benchmark & $N$ & PR obs. (calibrated) & PR / null & rms excess ($\times$ null) & latent dims \\
\midrule
ARC-Challenge & 1362 & 23.6 & 0.052 & $5.5\times$ & 6 \\
Winogrande & 1361 & 20.1 & 0.041 & $4.6\times$ & 5 \\
TruthfulQA & 1334 & 17.4 & 0.056 & $4.1\times$ & 5 \\
GSM8K & 1228 & 118.6 & 0.214 & $2.8\times$ & 4 \\
HellaSwag & 1362 & 27.3 & 0.028 & $11.8\times$ & 12 \\
\bottomrule
\end{tabular}
\caption{Participation ratio (PR) of the residual correlation spectrum after exact-margin
conditioning, against its own exact-margin null (mean of $400$ curveball replicates). ``Latent
dims'' counts eigenvalues exceeding the null's spectral edge (parallel analysis against the exact
null rather than iid noise). Before conditioning, the same statistic reads $1.5$--$3.6$ on every
benchmark, driven almost entirely by the shared item-difficulty eigenvalue.}
\label{tab:neff}
\end{table}

Table~\ref{tab:neff} makes the calibration gap concrete. The uncalibrated participation ratio
--- the same summary statistic, computed the naive way --- sits at $1.5$--$3.6$ on every benchmark:
read directly, that looks like near-total redundancy in the panel. After conditioning, real
structure appears: the calibrated ratio is only $2.8$--$21.4\%$ of its own null, meaning the panel
retains substantially more residual, non-difficulty-driven correlated structure than the calibrated
comparison would suggest by coincidence, and that structure is not one dominant shared factor but
$4$--$12$ separate latent dimensions (eigenvalues above the null's spectral edge). Neither number is
visible from the mean-level statistic alone.

\section{This is not just duplicate members}

A natural objection to any diversity shortfall is that it is driven by near-duplicate panel
members. We check this directly: clustering models at agreement $1.0000$ and keeping one
representative per cluster removes $900$ of ARC's $1{,}362$ models, and moves the calibrated ratio
from $0.052$ only to $0.058$. We also check whether ordinary discrimination heterogeneity among
otherwise-independent members is sufficient to produce the effect: a two-parameter item-response
simulation with realistic discrimination variance tops out at rms residual correlation $0.076$,
against an observed $0.248$ on ARC. What best matches the observed spectral signature
($\lambda_1^2/\sum\lambda^2=0.540$) is a simulated population with two \emph{additional} latent
ability dimensions ($0.457$), not a population of duplicated members ($0.102$). The calibrated
diagnostic is therefore picking up genuine multi-dimensional behavioral structure, not an artifact
of redundant submissions or of a single unmodeled discrimination parameter.

\section{Relevance to diversity-aware benchmark governance}

This workshop's scope explicitly includes diversity-aware and continuously evolving safety
benchmarks and coverage metrics for dynamic evaluation. The measurement failure mode we describe
is not specific to static multiple-choice accuracy: any coverage or diversity score computed from
raw pairwise (dis)agreement rates over a set of test cases --- attack variants, red-team probes,
behavioral test suites --- inherits the same degeneracy whenever test-case difficulty (or, in a
safety setting, attack potency) varies across the suite. A search procedure that appears to be
discovering diverse failure modes because raw disagreement is high may instead be finding items of
heterogeneous difficulty; a procedure that looks converged because raw agreement is high may be
sitting on a benchmark of uniform difficulty despite covering genuinely distinct behavioral axes.
Margin-calibration --- comparing against the exact conditional null implied by the test suite's own
per-item difficulty profile, rather than against an uninformative uniform baseline --- is a
domain-agnostic correction available as a diagnostic wherever a diversity or coverage score is built
from pairwise outcome comparisons, independent of whether the underlying search procedure is
evolutionary, adversarial, or manual.

\section{Limitations}

We study existing open models on static, single-turn multiple-choice benchmarks, not the
multi-turn, tool-using, or agentic settings central to this workshop; we have not applied this
calibration inside an active evolutionary or quality-diversity search loop, and doing so is the
natural next step rather than something we report here. Our conditioning model is unidimensional
(Rasch); Section 3's diagnostics rule out discrimination heterogeneity and duplicate members as
sufficient explanations but cannot fully separate genuine behavioral concentration from an
underfit null family. Claims are restricted to the evaluated open-model ecosystem. Code, data, and
the pre-registration are archived and linked in the full paper.

\begin{thebibliography}{9}
\small
\bibitem[Kuncheva \& Whitaker(2003)]{kuncheva2003} L.~I. Kuncheva and C.~J. Whitaker. Measures of diversity in classifier ensembles and their relationship with the ensemble accuracy. \emph{Machine Learning}, 51(2):181--207, 2003.
\bibitem[Strona et al.(2014)]{strona2014} G.~Strona, D.~Nappo, F.~Boccacci, S.~Fattorini, and J.~San-Miguel-Ayanz. A fast and unbiased procedure to randomize ecological binary matrices with fixed row and column totals. \emph{Nature Communications}, 5:4114, 2014.
\end{thebibliography}

\end{document}
